\documentclass{ximera}
%herschreven 13/11/26
\title{Beduidende cijfers en standaardvorm}
\author{Daan Lipkens}

\begin{document}
\begin{abstract}
    De nauwkeurigheid van meetresultaten, rekenen met beduidende cijfers, de wetenschappelijke notatie en de standaardvorm.
\end{abstract}
\maketitle

\section{Beduidende cijfers}\label{sec:beduidende_cijfers}
Elk meettoestel heeft een beperkte nauwkeurigheid. Met een meetlat met een millimeterverdeling lees je een lengte af tot op een millimeter, bijvoorbeeld $23{,}4 \, \mathrm{cm}$. Schrijf je $23{,}4000 \, \mathrm{cm}$, dan beweer je een nauwkeurigheid die je niet gemeten hebt. Het aantal cijfers dat je noteert, zegt dus iets over de nauwkeurigheid.

\begin{definition}[title={Beduidende cijfers}]\nl
    De \textbf{beduidende cijfers} (B.C.) van een meetresultaat zijn alle cijfers die je werkelijk afleest of berekent.
    \begin{itemize}
        \item Nullen \textbf{vooraan} tellen niet mee: ze geven enkel de plaats van de komma aan.
        \item Nullen \textbf{tussen} of \textbf{achter} andere cijfers tellen wel mee.
    \end{itemize}
\end{definition}

\begin{example}[title={Aantal beduidende cijfers}]\nl
    \begin{center}
        \begin{tabular}{|l|c|l|}
            \hline
            \textbf{Meetresultaat} & \textbf{Aantal B.C.} & \textbf{Uitleg} \\
            \hline
            $102 \, \mathrm{g}$ & 3 & de nul staat tussen andere cijfers \\
            $0{,}002 \, \mathrm{m}$ & 1 & nullen vooraan tellen niet \\
            $0{,}020 \, \mathrm{m}$ & 2 & de laatste nul staat achter de 2 \\
            $2{,}0 \, \mathrm{s}$ & 2 & de nul geeft de nauwkeurigheid aan \\
            $4{,}50 \cdot 10^{3} \, \mathrm{J}$ & 3 & de macht van 10 telt niet mee \\
            \hline
        \end{tabular}
    \end{center}
\end{example}

\begin{proposition}[title={Afronden}]\nl
    Is het eerste cijfer dat je weglaat kleiner dan 5, dan rond je naar beneden af. Is het 5 of groter, dan rond je naar boven af. Zo wordt $3{,}847$ afgerond op 2 B.C. gelijk aan $3{,}8$, en $3{,}85$ wordt $3{,}9$.
\end{proposition}

In een berekening combineer je meestal meerdere meetresultaten. De uitkomst kan nooit nauwkeuriger zijn dan het minst nauwkeurige gegeven.

\begin{proposition}[title={Som en verschil}]\nl
    Bij een \textbf{som} of een \textbf{verschil} rond je af op evenveel \textbf{cijfers na de komma} als het gegeven met het kleinste aantal cijfers na de komma. Zorg eerst dat alle termen dezelfde eenheid hebben.
\end{proposition}

\begin{example}[title={Som en verschil}]\nl
    \begin{align*}
        14{,}3 \, \mathrm{kg} + 3{,}56 \, \mathrm{kg} &= 17{,}86 \, \mathrm{kg} \simeq 17{,}9 \, \mathrm{kg} && \text{($14{,}3$ heeft 1 cijfer na de komma)}\\
        200{,}9 \, \mathrm{m} - 15{,}86 \, \mathrm{m} &= 185{,}04 \, \mathrm{m} \simeq 185{,}0 \, \mathrm{m} && \text{($200{,}9$ heeft 1 cijfer na de komma)}
    \end{align*}
\end{example}

\begin{proposition}[title={Product en quotiënt}]\nl
    Bij een \textbf{product} of een \textbf{quotiënt} rond je af op evenveel \textbf{beduidende cijfers} als het gegeven met het kleinste aantal beduidende cijfers.
\end{proposition}

\begin{example}[title={Product en quotiënt}]\nl
    \begin{align*}
        0{,}25 \, \mathrm{m} \cdot 0{,}334 \, \mathrm{m} &= 0{,}0835 \, \mathrm{m}^{2} \simeq 0{,}084 \, \mathrm{m}^{2} && \text{($0{,}25$ heeft 2 B.C.)}\\
        \frac{125 \, \mathrm{kg}}{44 \, \mathrm{kg}} &= 2{,}84 \simeq 2{,}8 && \text{($44$ heeft 2 B.C., de eenheden vallen weg)}
    \end{align*}
\end{example}

\begin{remark}[title={Exacte getallen}]\nl
    Getallen die niet gemeten zijn, zoals de $\frac{1}{2}$ in $E_{\text{kin}} = \frac{1}{2} \cdot m \cdot v^{2}$ of het aantal lampjes in een schakeling, zijn exact. Ze beperken het aantal beduidende cijfers niet.
\end{remark}

\begin{question}
    Hoeveel beduidende cijfers heeft $0{,}0450 \, \mathrm{m}$?

    \begin{multipleChoice}
        \choice{2}
        \choice[correct]{3}
        \choice{4}
        \choice{5}
    \end{multipleChoice}

    \begin{solution}\nl
        De twee nullen vooraan tellen niet mee. De cijfers 4, 5 en de laatste 0 tellen wel: 3 B.C.
    \end{solution}
\end{question}

\section{Wetenschappelijke notatie}\label{sec:wetenschappelijke_notatie}
Heel grote en heel kleine getallen schrijf je compact met een macht van 10. Die schrijfwijze toont meteen ook het aantal beduidende cijfers.

\begin{definition}[title={Wetenschappelijke notatie}]\nl
    Een getal staat in \textbf{wetenschappelijke notatie} als het geschreven is als
    \[ a \cdot 10^{n} \qquad \text{met } 1 \leq a < 10 \text{ en } n \text{ een geheel getal.} \]
    Er staat dus precies één cijfer, verschillend van nul, voor de komma. Alle cijfers van $a$ zijn beduidende cijfers.
\end{definition}

\begin{example}[title={Wel of niet in wetenschappelijke notatie?}]\nl
    \begin{itemize}
        \item $6{,}33 \cdot 10^{5} \, \mathrm{m}$ staat in wetenschappelijke notatie.
        \item $25{,}3 \cdot 10^{2} \, \mathrm{Pa}$ niet: er staan twee cijfers voor de komma. Correct is $2{,}53 \cdot 10^{3} \, \mathrm{Pa}$.
        \item $0{,}8 \cdot 10^{-3} \, \mathrm{C}$ niet: het cijfer voor de komma is nul. Correct is $8 \cdot 10^{-4} \, \mathrm{C}$.
    \end{itemize}
\end{example}

\begin{example}[foldable=true,title={Voorbeeld: afronden en noteren in wetenschappelijke notatie}]\nl
    Je meet $0{,}03589 \, \mathrm{m}^{2}$. Noteer dit resultaat in wetenschappelijke notatie met 2 B.C.

    \begin{solution}\nl
        \begin{align*}
            0{,}03589 \, \mathrm{m}^{2} &= 3{,}589 \cdot 10^{-2} \, \mathrm{m}^{2} && \text{(komma twee plaatsen naar rechts, dus $10^{-2}$)}\\
            &\simeq \boxed{3{,}6 \cdot 10^{-2} \, \mathrm{m}^{2}} && \text{(afronden op 2 B.C.)}
        \end{align*}
    \end{solution}
\end{example}

\begin{remark}[title={Nauwkeurigheid tonen}]\nl
    Schrijf je $4500 \, \mathrm{m}$, dan zie je niet of de nullen gemeten zijn. In wetenschappelijke notatie is dat wel duidelijk: $4{,}5 \cdot 10^{3} \, \mathrm{m}$ heeft 2 B.C., $4{,}500 \cdot 10^{3} \, \mathrm{m}$ heeft er 4.
\end{remark}

\section{Standaardvorm}\label{sec:standaardvorm}
\begin{definition}[title={Standaardvorm}]\nl
    Een resultaat staat in \textbf{standaardvorm} als het
    \begin{enumerate}
        \item het juiste aantal beduidende cijfers heeft,
        \item in wetenschappelijke notatie staat, en
        \item uitgedrukt is in de SI-eenheid.
    \end{enumerate}
\end{definition}

De standaardvorm heeft drie voordelen: je antwoord staat in de afgesproken eenheid, je leest de nauwkeurigheid meteen af en je kan het resultaat snel schatten door enkel de getallen voor de macht van 10 te vermenigvuldigen. In deze cursus staat elk eindantwoord in standaardvorm.

\begin{example}[foldable=true,title={Voorbeeld: een oppervlakte in standaardvorm}]\nl
    Een terrein is $24{,}23 \, \mathrm{m}$ lang en $43{,}1 \, \mathrm{m}$ breed. Bereken de oppervlakte en noteer ze in standaardvorm.

    \begin{solution}\nl
        \underline{Schatting:} $24{,}23 \, \mathrm{m} \cdot 43{,}1 \, \mathrm{m} \approx 2{,}4 \cdot 10^{1} \, \mathrm{m} \cdot 4{,}3 \cdot 10^{1} \, \mathrm{m} \approx 10 \cdot 10^{2} \, \mathrm{m}^{2} = 10^{3} \, \mathrm{m}^{2}$.

        \underline{Berekening:}
        \begin{align*}
            A &= 24{,}23 \, \mathrm{m} \cdot 43{,}1 \, \mathrm{m} = 1044{,}3 \, \mathrm{m}^{2} && \text{(product)}\\
            &\simeq \boxed{1{,}04 \cdot 10^{3} \, \mathrm{m}^{2}} && \text{($43{,}1$ heeft 3 B.C.)}
        \end{align*}
        Het resultaat ligt in de buurt van de schatting.
    \end{solution}
\end{example}

\begin{question}
    Welk resultaat staat in standaardvorm?

    \begin{multipleChoice}
        \choice{$1{,}50 \, \mathrm{mm}$}
        \choice{$0{,}00150 \, \mathrm{m}$}
        \choice{$15{,}0 \cdot 10^{-4} \, \mathrm{m}$}
        \choice[correct]{$1{,}50 \cdot 10^{-3} \, \mathrm{m}$}
    \end{multipleChoice}

    \begin{solution}\nl
        Alle vier stellen dezelfde lengte met 3 B.C. voor. Enkel $1{,}50 \cdot 10^{-3} \, \mathrm{m}$ staat in de SI-eenheid (meter) én in wetenschappelijke notatie.
    \end{solution}
\end{question}

\end{document}
